% !TEX root = ../main.tex
% Agradecimientos. Van en una página propia, después de los índices.
% En main.tex, justo después del abstract en inglés:
%     % !TEX root = ../main.tex
% Agradecimientos. Van en una página propia, después de los índices.
% En main.tex, justo después del abstract en inglés:
%     % !TEX root = ../main.tex
% Agradecimientos. Van en una página propia, después de los índices.
% En main.tex, justo después del abstract en inglés:
%     % !TEX root = ../main.tex
% Agradecimientos. Van en una página propia, después de los índices.
% En main.tex, justo después del abstract en inglés:
%     \input{preliminares/agradecimientos}
% El título es un \chapter*, con el mismo formato que los índices y resúmenes.

\chapter*{Agradecimientos}
\markboth{Agradecimientos}{Agradecimientos}
\begingroup

A mis tutores, David Gimeno y Salvador España, por aguantarme, por orientarme durante todo
el trabajo y por su dedicación constante.

A mis padres, que con su sacrificio me han dado la oportunidad de hacer todo esto, y a mi
hermano, por ser un pilar durante toda mi vida.

A mi amiga Carla, por acompañarme y hacer el proceso más ameno.

% Sugerencia (opcional): en las BCI es costumbre agradecer al participante.
% Al participante T15 y a su familia, y a los investigadores que han hecho públicos los datos
% de \textit{Brain-to-Text '25}, sin los que este trabajo no habría sido posible.

% ==================== BROMA: BORRAR ANTES DE ENTREGAR ====================

Al cerebro humano, por producir señales eléctricas.

A William Morton, Crawford Long, Horace Wells y Charles Jackson, por descubrir la anestesia..

Al equipo de marketing de Turkish Airlines, por motivos que no procede detallar aquí.

Al Decreto-ley 5/1968, de 6 de junio, por fundar la UPV (técnicamente fundó el Instituto
Politécnico Superior de Valencia, pero no vamos a ponernos tiquismiquis), y a Rafael Couchoud,
su primer rector, por hacerse cargo del invento.

A Jaume I, por conquistar Valencia en 1238 y unirla a la Corona de Aragón.

A Frodo, por destruir el Anillo Único.

A la Sociedad para la Prevención de Pájaros, cuya labor ha sido impecable: en la elaboración
de este trabajo no ha participado ni un solo pájaro.

Y, por último, como no podía ser de otra manera, a Saturno y, obviamente, a todos sus anillos
% ======================== FIN DE LA BROMA ================================
\endgroup
\cleardoublepage

% El título es un \chapter*, con el mismo formato que los índices y resúmenes.

\chapter*{Agradecimientos}
\markboth{Agradecimientos}{Agradecimientos}
\begingroup

A mis tutores, David Gimeno y Salvador España, por aguantarme, por orientarme durante todo
el trabajo y por su dedicación constante.

A mis padres, que con su sacrificio me han dado la oportunidad de hacer todo esto, y a mi
hermano, por ser un pilar durante toda mi vida.

A mi amiga Carla, por acompañarme y hacer el proceso más ameno.

% Sugerencia (opcional): en las BCI es costumbre agradecer al participante.
% Al participante T15 y a su familia, y a los investigadores que han hecho públicos los datos
% de \textit{Brain-to-Text '25}, sin los que este trabajo no habría sido posible.

% ==================== BROMA: BORRAR ANTES DE ENTREGAR ====================

Al cerebro humano, por producir señales eléctricas.

A William Morton, Crawford Long, Horace Wells y Charles Jackson, por descubrir la anestesia..

Al equipo de marketing de Turkish Airlines, por motivos que no procede detallar aquí.

Al Decreto-ley 5/1968, de 6 de junio, por fundar la UPV (técnicamente fundó el Instituto
Politécnico Superior de Valencia, pero no vamos a ponernos tiquismiquis), y a Rafael Couchoud,
su primer rector, por hacerse cargo del invento.

A Jaume I, por conquistar Valencia en 1238 y unirla a la Corona de Aragón.

A Frodo, por destruir el Anillo Único.

A la Sociedad para la Prevención de Pájaros, cuya labor ha sido impecable: en la elaboración
de este trabajo no ha participado ni un solo pájaro.

Y, por último, como no podía ser de otra manera, a Saturno y, obviamente, a todos sus anillos
% ======================== FIN DE LA BROMA ================================
\endgroup
\cleardoublepage

% El título es un \chapter*, con el mismo formato que los índices y resúmenes.

\chapter*{Agradecimientos}
\markboth{Agradecimientos}{Agradecimientos}
\begingroup

A mis tutores, David Gimeno y Salvador España, por aguantarme, por orientarme durante todo
el trabajo y por su dedicación constante.

A mis padres, que con su sacrificio me han dado la oportunidad de hacer todo esto, y a mi
hermano, por ser un pilar durante toda mi vida.

A mi amiga Carla, por acompañarme y hacer el proceso más ameno.

% Sugerencia (opcional): en las BCI es costumbre agradecer al participante.
% Al participante T15 y a su familia, y a los investigadores que han hecho públicos los datos
% de \textit{Brain-to-Text '25}, sin los que este trabajo no habría sido posible.

% ==================== BROMA: BORRAR ANTES DE ENTREGAR ====================

Al cerebro humano, por producir señales eléctricas.

A William Morton, Crawford Long, Horace Wells y Charles Jackson, por descubrir la anestesia..

Al equipo de marketing de Turkish Airlines, por motivos que no procede detallar aquí.

Al Decreto-ley 5/1968, de 6 de junio, por fundar la UPV (técnicamente fundó el Instituto
Politécnico Superior de Valencia, pero no vamos a ponernos tiquismiquis), y a Rafael Couchoud,
su primer rector, por hacerse cargo del invento.

A Jaume I, por conquistar Valencia en 1238 y unirla a la Corona de Aragón.

A Frodo, por destruir el Anillo Único.

A la Sociedad para la Prevención de Pájaros, cuya labor ha sido impecable: en la elaboración
de este trabajo no ha participado ni un solo pájaro.

Y, por último, como no podía ser de otra manera, a Saturno y, obviamente, a todos sus anillos
% ======================== FIN DE LA BROMA ================================
\endgroup
\cleardoublepage

% El título es un \chapter*, con el mismo formato que los índices y resúmenes.

\chapter*{Agradecimientos}
\markboth{Agradecimientos}{Agradecimientos}
\begingroup

A mis tutores, David Gimeno y Salvador España, por aguantarme, por orientarme durante todo
el trabajo y por su dedicación constante.

A mis padres, que con su sacrificio me han dado la oportunidad de hacer todo esto, y a mi
hermano, por ser un pilar durante toda mi vida.

A mi amiga Carla, por acompañarme y hacer el proceso más ameno.

% Sugerencia (opcional): en las BCI es costumbre agradecer al participante.
% Al participante T15 y a su familia, y a los investigadores que han hecho públicos los datos
% de \textit{Brain-to-Text '25}, sin los que este trabajo no habría sido posible.

% ==================== BROMA: BORRAR ANTES DE ENTREGAR ====================

Al cerebro humano, por producir señales eléctricas.

A William Morton, Crawford Long, Horace Wells y Charles Jackson, por descubrir la anestesia..

Al equipo de marketing de Turkish Airlines, por motivos que no procede detallar aquí.

Al Decreto-ley 5/1968, de 6 de junio, por fundar la UPV (técnicamente fundó el Instituto
Politécnico Superior de Valencia, pero no vamos a ponernos tiquismiquis), y a Rafael Couchoud,
su primer rector, por hacerse cargo del invento.

A Jaume I, por conquistar Valencia en 1238 y unirla a la Corona de Aragón.

A Frodo, por destruir el Anillo Único.

A la Sociedad para la Prevención de Pájaros, cuya labor ha sido impecable: en la elaboración
de este trabajo no ha participado ni un solo pájaro.

Y, por último, como no podía ser de otra manera, a Saturno y, obviamente, a todos sus anillos
% ======================== FIN DE LA BROMA ================================
\endgroup
\cleardoublepage
